\documentclass[10pt,a4paper,twocolumn]{article}
\usepackage[margin=2cm,columnsep=0.7cm]{geometry}

\usepackage[T1]{fontenc}
\usepackage[utf8]{inputenc}
\usepackage{amsmath,amssymb,amsthm}
\usepackage{graphicx}
\usepackage{booktabs}
\usepackage{xcolor}
\usepackage[hidelinks]{hyperref}
\usepackage[numbers]{natbib}

\newcommand{\todo}[1]{\textcolor{red}{[#1]}}

\title{When does the latent space matter for pseudotime? A benchmark with clonal ground truth}
\author{Vladimir Opait}
\date{\today}

\begin{document}

\twocolumn[{
\begin{@twocolumnfalse}
\maketitle

\begin{abstract}
\todo{Draft of the new framing, numbers from the dev split.}
Trajectory inference methods order cells along a low-dimensional representation of the
expression matrix, but the choice of this representation is rarely evaluated. We use the clonal
barcodes of the LARRY in vitro haematopoiesis data to score pseudotime by the time order of
sister cells, and ask which latent space helps, where and why. We compare [N] latent spaces,
from PCA to count-based variational autoencoders, graph autoencoders, diffusion smoothing,
two-dimensional layouts and a single-cell foundation model, as inputs to Euclidean distance, DPT
and Palantir, at full sequencing depth and after binomial thinning of the counts. At full depth
the trajectory method matters more than the latent space: with DPT and Palantir no latent space
improves on PCA by more than 0.015 in accuracy, which is comparable to the spread over roots and training seeds.
When molecules are removed, the latent space starts to matter: at 20\% of the molecules scVI
improves DPT by 0.09 over PCA in each of three thinning draws, while diffusion smoothing does not
help. [Which part of scVI is responsible.] [Foundation model result.] We conclude that
[recommendation].
\end{abstract}
\vspace{1.5em}
\end{@twocolumnfalse}
}]

\section{Introduction}
\label{sec:intro}
\todo{Narrative introduction. (1) Trajectory inference orders cells along a low-dimensional
representation of the expression matrix. (2) The number of available representations grows: PCA,
count-based variational autoencoders, graph models, diffusion smoothing, two-dimensional layouts,
single-cell foundation models. (3) They are usually evaluated on simulations or on annotations,
without a ground truth for the order of real cells; zero-shot foundation model embeddings were
found worse than HVG--PCA for OT-based dynamics \cite{TODO-zhou2026}. (4) Clonal barcodes of
LARRY give such a ground truth: a sister cell sampled later should have larger pseudotime.
(5) Our question: when does the latent space matter for pseudotime, and what kind of latent space
helps. (6) Short summary of the answer.}

\section{Results}
\todo{All numbers are from the development clones (1{,}760 clones, 26{,}106 pairs); test not opened.}

\subsection{A clonal benchmark for latent spaces}
\label{sec:overview}
\todo{Short overview: what is compared, the sister-pair metric, the grid; details in
Methods. Order inside each subsection: what was done, observations with numbers, one-sentence
conclusion, explanation.}

\begin{figure*}[t]
\centering
\todo{Fig.~1: (a) sister pair from a clone sampled at two days, correct if the later cell has
larger pseudotime; (b) pipeline: counts $\to$ latent space $\to$ trajectory method $\to$ metric;
(c) thinning of counts.}
\caption{Benchmark design.}
\label{fig:design}
\end{figure*}

Table~\ref{tab:grid} lists the settings that are varied. Trained latent spaces are trained with 3 seeds at full depth. Since spread
between roots resulted to be small, we don't rerun model with different roots. A trained latent space without
an internal graph does not depend on $k$, so one model serves all values of $k$.

\begin{table*}[t]
\centering\small
\caption{Benchmark grid.}
\label{tab:grid}
\begin{tabular}{lp{0.62\linewidth}}
\toprule
Setting & Values \\
\midrule
Graph size $k$      & 15 (10, 30 in the appendix) \\
Latent space, main  & PCA-50; PCA-10; PCA of Pearson residuals; LinearSCVI; scVI; MAGIC;
                      GAE-50; Geneformer (zero-shot, LoRA) \\
Latent space, full depth only & PCA-5, PCA-20; smoothed PCA ($t = 1, 3, 10$); scVI variants;
                      scTour; GAE-10; VGAE-10, VGAE-50; untrained GAE \\
Latent space, 2D    & UMAP; \todo{PCA-2; SPRING} (with their own set of trajectory methods,
                      Section~\ref{sec:twod}) \\
Sequencing depth $p$ & 1; 0.5, 0.3, 0.2 (3 draws each); 0.7, 0.1 in the appendix \\
Trajectory method   & Euclidean distance; DPT; Palantir; \todo{Slingshot; Waddington-OT} \\
Root                & 5 fixed cells \\
Training seed       & 0--2 at full depth (scVI, GAE-50); 0 under thinning \\
\bottomrule
\end{tabular}
\end{table*}


\subsection{At full depth the trajectory method matters more than the latent space}
On PCA, Euclidean distance reaches acc 0.826, DPT 0.876 and Palantir 0.888
(Table~\ref{tab:full}). With DPT and Palantir, no latent space improves on PCA by more than
0.015. The largest gain is GAE-50 with DPT ($+0.014$, interval $+0.008$ to $+0.020$), which is
reproduced with three training seeds ($+0.013$ to $+0.015$). Euclidean distance is the only method
where the latent space changes acc substantially: fewer PCA components or smoothing gain up to
0.04, and graph autoencoders lose about 0.2, because they
are trained to preserve neighbourhoods, not distances between distant cells. scTour, the only
latent space made for trajectories, is worse than PCA with DPT ($-0.055$) and Palantir ($-0.063$).
\todo{MAGIC, Pearson residuals, LinearSCVI, Geneformer.}

\begin{table*}[t]
\centering\small
\caption{Accuracy at full depth ($k=15$, mean over five roots). PCA: acc with 95\% interval;
other rows: paired difference to PCA with 95\% clone bootstrap interval. All latent spaces in
Table~\ref{tab:fullall}.}
\label{tab:full}
\scriptsize
\setlength{\tabcolsep}{3pt}
\begin{tabular}{llll}
\toprule
Latent space & Euclidean distance & DPT & Palantir \\
\midrule
PCA-50 & 0.826 (0.795, 0.861) & 0.876 (0.846, 0.912) & 0.888 (0.860, 0.920) \\
\midrule
PCA-10 & $+$0.026 ($+$0.019, $+$0.036) & $-$0.002 ($-$0.006, $+$0.002) & $-$0.019 ($-$0.026, $-$0.014) \\
PCA, Pearson residuals & \todo{} & \todo{} & \todo{} \\
LinearSCVI & \todo{} & \todo{} & \todo{} \\
scVI & $-$0.024 ($-$0.039, $-$0.011) & $-$0.004 ($-$0.011, $+$0.004) & $-$0.021 ($-$0.029, $-$0.013) \\
MAGIC & \todo{} & \todo{} & \todo{} \\
GAE-50 & $-$0.212 ($-$0.267, $-$0.164) & $+$0.014 ($+$0.008, $+$0.020) & $-$0.000 ($-$0.003, $+$0.004) \\
Geneformer, zero-shot & \todo{} & \todo{} & \todo{} \\
Geneformer, LoRA & \todo{} & \todo{} & \todo{} \\
\bottomrule
\end{tabular}
\end{table*}

\subsection{Latent effects at full depth are of the size of root, seed and thinning noise}
\todo{add a short summary of the design checks: settings of the trajectory methods, $k$, DPT kernel (Appendix~\ref{sec:checks}).}
Differences between latent spaces at full depth are of the same size as other sources of
variation that are usually not reported. The standard deviation of acc over the five roots is
0.008--0.012 on PCA (Table~\ref{tab:roots}). Retraining scVI with another seed changes acc by up
to 0.012 with DPT and 0.029 with Euclidean distance; GAE-50 changes by at most 0.001 with DPT and
Palantir (Table~\ref{tab:seeds}). Two independent thinning draws at the same depth give acc that
differs by up to 0.02 for $p \ge 0.3$, and up to 0.13 for $p = 0.2$ (Palantir on smoothed PCA,
$t=3$: 0.799, 0.727, 0.665). Paired differences to PCA on the same draw are much more stable than
acc itself. \todo{variance decomposition: method, latent, root, draw, seed}

\subsection{At low depth a count model helps and smoothing does not}
With fewer molecules, acc on PCA falls for all methods: for DPT from 0.876 at full depth to 0.826
at $p=0.3$ and 0.763 at $p=0.2$ (Table~\ref{tab:thin}, Fig.~\ref{fig:thin}). scVI loses less. With
DPT it is better than PCA by $+0.049$ at $p=0.3$ and $+0.094$ at $p=0.2$, and the gain is positive
in each of the three draws ($+0.087$ to $+0.098$ at $p=0.2$). With Palantir the gain is $+0.030$
and $+0.065$. With Euclidean distance scVI is worse than PCA at every depth. Diffusion smoothing of
PCA does not help DPT or Palantir at low depth and costs up to 0.05 with Palantir at $p=0.2$
(Table~\ref{tab:thinall}). GAE-50 loses its small advantage already at $p=0.3$ (one draw).
\todo{MAGIC, Pearson residuals, LinearSCVI, Geneformer.}

To see whether the loss comes from cells with few molecules, we split the sister pairs into ten
equal groups by the number of molecules in the shallower cell of the pair and followed the same
groups across depths (Fig.~\ref{fig:depthgroups}). Within one depth, accuracy hardly depends on
the group: with DPT on PCA at full depth it is 0.86--0.92 in groups 1 to 9. Only the deepest group
is worse (0.76); it contains more pairs in which both cells are undifferentiated (31\% against
17--25\% in the other groups), and such pairs are hard to order. When molecules are removed,
accuracy falls in all groups at once: at $p=0.1$ it is 0.49--0.61 for PCA in every group. The loss
therefore does not come from the shallow cells themselves but from the depth of the whole data
set, which is consistent with a degraded neighbourhood graph \todo{not checked directly}. scVI
keeps the same profile at $p=0.3$ as at full depth and stays at 0.69--0.84 at $p=0.1$.

\begin{figure*}[t]
\centering
\includegraphics[width=\linewidth]{../pics/fig_depth_acc.pdf}
\caption{Accuracy of DPT by groups of sister pairs. Dev pairs are split into ten equal groups by
the number of molecules in the shallower cell of the pair (group 1: fewest, group 10: most; at
full depth the group medians range from 691 to 5{,}862 molecules); the same pairs form each group
at every depth. Lines: kept fraction of molecules $p$ (thinning draw 0); mean over five roots,
95\% clone bootstrap band. The lower accuracy of group 10 reflects pairs of two undifferentiated
cells, not depth. All trajectory methods and latent spaces in Fig.~\ref{fig:depthall}.}
\label{fig:depthgroups}
\end{figure*}

\begin{table*}[t]
\centering\small
\caption{Accuracy after thinning ($k=15$). PCA: acc; other rows: paired difference to PCA on the
same draw. Mean over three draws, with minimum and maximum in brackets. $^\dagger$One draw.
$p=0.7$ and $p=0.1$ in Table~\ref{tab:thinall}.}
\label{tab:thin}
\footnotesize
\setlength{\tabcolsep}{3pt}
\begin{tabular}{lllll}
\toprule
Latent space & $p$ & Euclidean distance & DPT & Palantir \\
\midrule
PCA-50 & 1   & 0.826 & 0.876 & 0.888 \\
       & 0.5 & 0.810 [0.802, 0.814] & 0.873 [0.869, 0.876] & 0.874 [0.870, 0.878] \\
       & 0.3 & 0.779 [0.775, 0.782] & 0.826 [0.821, 0.833] & 0.845 [0.841, 0.851] \\
       & 0.2 & 0.768 [0.763, 0.774] & 0.763 [0.755, 0.774] & 0.780 [0.766, 0.795] \\
\midrule
scVI & 1   & $-$0.024 & $-$0.004 & $-$0.021 \\
     & 0.5 & $-$0.047 [$-$0.056, $-$0.041] & $+$0.004 [$+$0.002, $+$0.006] & $-$0.017 [$-$0.024, $-$0.014] \\
     & 0.3 & $-$0.021 [$-$0.034, $-$0.010] & $+$0.049 [$+$0.034, $+$0.064] & $+$0.030 [$+$0.018, $+$0.039] \\
     & 0.2 & $-$0.041 [$-$0.051, $-$0.030] & $+$0.094 [$+$0.087, $+$0.098] & $+$0.065 [$+$0.025, $+$0.091] \\
\midrule
GAE-50$^\dagger$ & 1   & $-$0.212 & $+$0.014 & $-$0.000 \\
                 & 0.5 & $-$0.234 & $-$0.001 & $-$0.004 \\
                 & 0.3 & $-$0.167 & $-$0.033 & $-$0.066 \\
                 & 0.2 & $-$0.177 & $-$0.116 & $-$0.207 \\
\midrule
\todo{other latents} & & & & \\
\bottomrule
\end{tabular}
\end{table*}

\begin{figure*}[t]
\centering
\todo{figure: acc against $p$ for each latent space, one panel per trajectory method; points are
draws, line is the mean}
\caption{Accuracy against sequencing depth.}
\label{fig:thin}
\end{figure*}

\subsection{What makes scVI robust to depth}
\todo{title as a conclusion once the ablation is done}
\todo{scVI differs from PCA-50 in three ways: dimension (10 vs 50), noise model (counts vs
log-normalized), non-linearity. PCA-10 tests the dimension, PCA of Pearson residuals and
LinearSCVI test the noise model without non-linearity.}

\subsection{Two-dimensional layouts}
\label{sec:twod}
Trajectories are often inferred or displayed in a two-dimensional layout. UMAP gives the best
Euclidean distance of all latent spaces: acc 0.876, $+0.050$ over PCA, with the same gain for
three seeds (Table~\ref{tab:twod}). Possibly UMAP separates cell states, so that the distance from
the root grows in steps between states and most sister pairs, whose later cell is in a more
mature state, are ordered correctly \todo{not checked}. With graph-based methods UMAP is worse
than PCA: $-0.095$ with DPT and $-0.019$ to $-0.032$ with Palantir. In two dimensions each cell has
neighbours only nearby in the plane, diffusion on the graph mixes slowly, and the second
eigenvalue of the Palantir diffusion operator is 0.99998 (double precision). In single precision it
became 1 for two of the three seeds and gave infinite coordinates in the multiscale space; we
recomputed the eigendecomposition in double precision in such cases. Palantir is designed to run on PCA, with UMAP used only for
display \cite{setty2019palantir}, so its result on UMAP shows the cost of the misuse rather than a
property of the latent space.

UMAP is less robust to depth than PCA. Its gain with Euclidean distance holds at $p=0.5$
($+0.057$) and becomes a loss at $p=0.2$ ($-0.099$); with DPT and Palantir the loss grows to
$-0.12$ and $-0.09$ at $p=0.2$. The spread between draws is also larger: with DPT at $p=0.3$ acc
ranges from 0.645 to 0.700, against 0.821 to 0.833 for PCA. \todo{PCA-2 as linear 2D control;
SPRING of the original study; Slingshot and Monocle 3 \cite{TODO-cao2019monocle3}, which are
designed for low-dimensional spaces.}

\begin{table*}[t]
\centering\small
\caption{Two-dimensional layouts ($k=15$, five roots). Paired difference to PCA-50 with the same
trajectory method. Full depth: seeds 0--2, with 95\% clone bootstrap interval for seed 0 (the
intervals of seeds 1 and 2 differ by at most 0.01). Thinning: seed 0,
mean over three draws with minimum and maximum ($p=0.7$: one draw).}
\label{tab:twod}
\scriptsize
\setlength{\tabcolsep}{3pt}
\begin{tabular}{lllll}
\toprule
Latent space & $p$ & Euclidean distance & DPT & Palantir \\
\midrule
UMAP, seed 0 & 1 & $+$0.050 ($+$0.037, $+$0.066) & $-$0.096 ($-$0.138, $-$0.057) & $-$0.028 ($-$0.041, $-$0.017) \\
UMAP, seed 1 & 1 & $+$0.051 & $-$0.094 & $-$0.019 \\
UMAP, seed 2 & 1 & $+$0.050 & $-$0.096 & $-$0.032 \\
\midrule
UMAP, seed 0 & 0.7 & $+$0.050 & $-$0.105 & $-$0.027 \\
             & 0.5 & $+$0.057 [$+$0.054, $+$0.060] & $-$0.130 [$-$0.138, $-$0.124] & $-$0.008 [$-$0.018, $-$0.002] \\
             & 0.3 & $+$0.017 [$+$0.000, $+$0.044] & $-$0.146 [$-$0.179, $-$0.121] & $-$0.057 [$-$0.077, $-$0.027] \\
             & 0.2 & $-$0.099 [$-$0.151, $-$0.071] & $-$0.122 [$-$0.144, $-$0.109] & $-$0.093 [$-$0.121, $-$0.069] \\
             & 0.1 & $-$0.287 [$-$0.330, $-$0.242] & $-$0.038 [$-$0.083, $+$0.035] & $-$0.104 [$-$0.129, $-$0.089] \\
\midrule
\todo{PCA-2, SPRING} & & & & \\
\bottomrule
\end{tabular}
\end{table*}

\subsection{A foundation model}
\todo{Geneformer zero-shot vs LoRA; compare with \cite{TODO-zhou2026}, who found zero-shot
embeddings worse than HVG--PCA on the same data with OT-based metrics.}

\section{Discussion}
\todo{(1) at full depth the latent space is a minor choice compared with the trajectory method,
and its effect is of the size of root and seed noise; (2) at low depth a count model helps and
smoothing does not, because smoothing on the same graph adds no information; (3) report roots,
seeds and thinning draws; (4) limitations: one dataset, in vitro, dev split, batch effect in LK.}

\section{Methods}

\subsection{Dataset}
We use the in vitro hematopoietic differentiation time course of LARRY \cite{weinreb2020lineage}.
Mouse progenitors were barcoded with the LARRY lentiviral library and profiled by inDrops at
days 2, 4 and 6 of culture, starting from two populations: Lin$^-$Kit$^+$Sca1$^+$ (LSK) and
Lin$^-$Kit$^+$Sca1$^-$ (LK). The data contain 130{,}887 cells and 25{,}289 genes in 35 sequencing
libraries. We use the total-count normalized matrix provided by the authors. Cells carry one of
eleven type annotations; 72{,}193 are undifferentiated, and the largest mature types are
neutrophils (22{,}237) and monocytes (19{,}182).

Clonal barcodes are available for 49{,}302 cells in 5{,}864 clones (median size 4, maximum 177).
Of these clones, 3{,}894 have cells at more than one day, which gives 208{,}788 pairs of
sister cells sampled at different days. Cells of a clone are not synchronous in differentiation,
so a clone does not define a reference time for each cell, but it fixes which cell was sampled
earlier.

\paragraph{Batch structure.}
Each sequencing library corresponds to one day and one starting population, so batch correction
would also remove the signal of time and population; we therefore did not correct for batch and
measured its size instead (Appendix~\ref{sec:batch}). Among cells of the same population, day,
well and type, the share of nearest neighbours from the same library is 1.03 times the random
expectation for LSK and 1.45 for LK, and library explains 0.2\% and 4.6\% of the PCA variance.
The batch effect is negligible for LSK and moderate for LK, which we treat as a limitation.

\subsection{Evaluation}
The pseudotime of a cell is compared with the time order within its clone. For a pair of sister
cells from the same clone, sampled at different days and assigned to the same branch, the pair
is correct if the cell from the later day has the larger pseudotime (ties count as 0.5). The
accuracy \emph{acc} is the share of correct pairs. To measure how large the errors are, the
pseudotime is converted to percentile ranks within the branch, and \emph{violation} is the mean
of $\max(0, r_{\mathrm{early}} - r_{\mathrm{late}})$ over pairs. A random pseudotime gives
acc $=0.5$ and violation $=1/6$. Confidence intervals come from a bootstrap over clones, not over
cells or pairs.

Branches are defined by the trajectory files of the original study. The neutrophil branch
contains the 61{,}310 cells of their neutrophil trajectory (63{,}079 sister pairs); the monocyte
branch contains the 36{,}148 remaining cells of their neutrophil--monocyte trajectory
(34{,}089 pairs). Metrics are computed per branch and averaged with equal weight; a branch needs
at least 200 pairs to enter the mean. As a check of the metric, the original pseudotime of the
neutrophil trajectory gives acc $=0.940$ (95\% CI 0.924--0.956) and violation $=0.004$, and a
random pseudotime gives acc $=0.501$ and violation $=0.167$.

\subsection{Sequencing depth}
\label{sec:thin}
To vary the sequencing depth, we thin the counts. Raw counts are recovered from the normalized
matrix by dividing each cell by its smallest non-zero value (the deviation from integers is below
$10^{-3}$). Each molecule is kept independently with probability $p$, i.e.\ the count of every
gene in every cell is replaced by a binomial draw with $p$. The median raw total is 2{,}249
molecules per cell (minimum 300). For $p \le 0.2$ a cell would keep fewer than 100 molecules on
average if its raw total is below $100/p$; such cells are thinned with probability $100/n$
instead, where $n$ is the raw total of the cell (\todo{share of cells: 4\% at $p=0.2$}). Without
this floor, cells with no counts left in the selected genes become identical points and Palantir
fails. The thinned counts are normalized to the median total of the full data times $p$, and
every latent space is recomputed from them, including the selection of variable genes. Roots
are kept from the full data. We use $p \in \{0.7, 0.5, 0.3, 0.2, 0.1\}$ and three independent
draws for $p \le 0.5$. Latent spaces are compared with PCA on the same draw.

\todo{evaluation for Waddington-OT: decided (W1 for held-out day 4 after GPA alignment, as in Zhou 2026; see claude.md, Metrics, item 4), to be written up}

\subsection{Trajectory methods}
For each method we describe what it uses from the latent space. The assumptions for the latent
spaces in Section~\ref{sec:latent} refer to these properties.

\paragraph{Euclidean distance.}
The pseudotime of a cell is its Euclidean distance to the root in the latent space. No graph is
used. This baseline shows how much of the order comes from the latent space alone.
\todo{root as for DPT}

\emph{What it uses.} \todo{which properties of the latent space this baseline relies on}

\paragraph{DPT.}
Diffusion pseudotime (DPT) \cite{haghverdi2016diffusion} was computed with Scanpy on all cells.
A $k$-nearest-neighbour graph was built on the latent space under evaluation, and DPT
distance was computed from 10 diffusion components. Edge weights use the Scanpy default (UMAP)
kernel; the Gaussian kernel of the original method changes acc by at most 0.006
(Table~\ref{tab:kernel}).

\emph{What it uses.} \todo{which properties of the latent space DPT relies on}

\paragraph{Slingshot.}
\todo{description}

\emph{What it uses.} \todo{which properties of the latent space Slingshot relies on}

\paragraph{Palantir.}
Palantir \cite{setty2019palantir} computes pseudotime as the shortest-path distance from the root
on a $k$-nearest-neighbour graph built on diffusion components, and also gives fate probabilities.
\todo{parameters; how the latent space is passed in}
Palantir computes the diffusion eigenvalues in single precision with tolerance $10^{-4}$. If any of
the used eigenvalues is 1 or larger, its scaling $\lambda/(1-\lambda)$ in the multiscale space is
infinite or negative; in such cases we recompute the eigendecomposition of the same kernel in double
precision with tolerance $10^{-8}$. This happened only for two-dimensional latent spaces
(Section~\ref{sec:twod}). If the iterative eigensolver for the waypoint Markov chain does not
converge, a dense eigendecomposition is used.

\emph{What it uses.} \todo{which properties of the latent space Palantir relies on}

\paragraph{Waddington-OT.}
Waddington-OT \cite{schiebinger2019optimal} links cells of consecutive days by optimal transport.
The cost of moving a cell to another is the squared Euclidean distance between them in the
latent space, so the cost matrix is the only part that depends on the latent space. The
transport is entropy-regularized (parameter $\varepsilon$) and unbalanced (parameter $\tau$),
because cells divide and die between days. Growth rates are estimated from two KEGG pathways
for mouse \cite{kanehisa2000kegg}: cell cycle (mmu04110) for proliferation and apoptosis
(mmu04210) for death. The same two pathways were used to estimate growth on these data in
PRESCIENT \cite{yeo2021prescient}. Of the 157 cell cycle genes, 154 are present in the data, and
all 136 apoptosis genes are present (four genes appear under older symbols: Casc5, Sgol1, Skp1a,
Sept4). Five genes (Atm, Gadd45a, Gadd45b, Gadd45g, Trp53) belong to both pathways. The gene
lists were retrieved from KEGG on 6 November 2026. We compute the growth rates once and use the same values and the same $\varepsilon$ and
$\tau$ for all latent spaces. The problem is solved with moscot \cite{klein2025moscot}.
\todo{values of $\varepsilon$ and $\tau$; how scores are turned into growth rates}

\emph{What it uses.} 
WOT assumes Markov nature of the data, requieres enough density in time, so it wouldn't need to do too big 
jumps in space, and it strongly depends on right hyperparameters.

\subsection{Latent spaces}
\label{sec:latent}
For each latent space we state a general assumption and, where a trajectory method is expected to
react differently, an assumption for that method.

\paragraph{PCA.}
Normalized counts were log-transformed ($\log(1+x)$), the 2000 most variable genes were
selected with the Scanpy default (Seurat) criterion \cite{wolf2018scanpy}, each gene was scaled
to unit variance with values clipped at 10, and the first 50 principal components were
computed by exact eigendecomposition of the covariance matrix. All 130{,}887 cells were used.

\emph{Assumption.} \todo{what this latent space changes in the geometry and why it should or should not help}

\emph{By method.} \todo{only where a method is expected to react differently}

\paragraph{Smoothed PCA.}
A $k$-nearest-neighbour graph was built on the PCA coordinates by exact Euclidean
search, with Scanpy connectivity weights $C$. The latent space was smoothed by $t$ steps of
diffusion, $X_t = P^t X$ with $P = D^{-1}(C+I)$, where $D$ is the diagonal matrix of row sums
of $C+I$. We use $t \in \{1, 3, 10\}$. This is a graph filter without learned weights and serves
as the control for graph autoencoders.

\emph{Assumption.} \todo{what this latent space changes in the geometry and why it should or should not help}

\emph{By method.} \todo{only where a method is expected to react differently}

\paragraph{MAGIC.}
MAGIC \cite{vandijk2018magic} imputes expression by diffusion on a cell--cell affinity graph.
\todo{parameters and how the latent is taken from the output}

\emph{Assumption.} \todo{what this latent space changes in the geometry and why it should or should not help}

\emph{By method.} \todo{only where a method is expected to react differently}

\paragraph{PCA with fewer components.}
The first 5, 10 or 20 components of the PCA above. This separates the effect of the
dimension from the effect of the model in the other latent spaces, most of which have 10
dimensions.

\paragraph{PCA of Pearson residuals.}
\todo{analytic Pearson residuals \cite{TODO-lause2021} of raw counts, 2000 genes, 50 components;
linear, but with a count noise model}

\paragraph{scVI.}
scVI \cite{TODO-lopez2018} is a variational autoencoder with a negative binomial likelihood on
raw counts. It was trained on the raw counts of the 2000 variable genes with the default
settings of scvi-tools (10 latent dimensions); the latent space is the posterior mean. Variants
with 10 or 30 dimensions and a maximum KL weight of 0.01 or 1 are in the appendix.

\paragraph{LinearSCVI.}
\todo{scVI with a linear decoder \cite{TODO-svensson2020}; same likelihood, no non-linearity}

\paragraph{scTour.}
scTour \cite{TODO-li2023sctour} is a variational autoencoder with a neural ODE for time, made for
trajectory inference. Default settings; the latent space is the mix of the encoder and ODE
latents with weights 0.5. \todo{genes, epochs}

\paragraph{Graph autoencoders.}
A graph autoencoder (GAE) and its variational version (VGAE) \cite{TODO-kipf2016} with two graph
convolution layers (hidden size 128) take standardized PCA coordinates and the $k$-nearest-neighbour
graph of PCA, and are trained for 200 epochs (Adam, learning rate 0.01) to reconstruct the edges of
the graph against random negative edges. Latent dimension 10 or 50. An untrained GAE (0 epochs)
is a random graph filter and serves as a control.

\paragraph{UMAP.}
UMAP \cite{TODO-mcinnes2018} was computed with Scanpy from the $k=15$ nearest-neighbour graph of
the PCA above (2 dimensions, \texttt{min\_dist} 0.5, Scanpy defaults otherwise), with seeds 0--2
at full depth and seed 0 under thinning. UMAP is the default space of Monocle 3
\cite{TODO-cao2019monocle3} and is widely used to display trajectories.

\paragraph{SPRING.}
\todo{two-dimensional force-directed layout of the original study \cite{TODO-weinreb2018spring},
taken from the metadata; full depth only}

\paragraph{Geneformer.}
\todo{Geneformer V2-104M \cite{TODO-theodoris2023}, as in \cite{TODO-zhou2026}: raw counts of
14{,}787 genes with one-to-one human orthologs (Ensembl BioMart; 83\% of molecules per cell, median),
rank-value tokenization, [CLS] embedding, PCA to 50 components. Zero-shot and after LoRA
\cite{TODO-hu2021lora} with the masked-gene objective on LARRY.}

\bibliographystyle{plain}
\bibliography{references}

\appendix
\renewcommand{\thetable}{S\arabic{table}}
\renewcommand{\thefigure}{S\arabic{figure}}
\setcounter{table}{0}
\setcounter{figure}{0}

\section{Batch structure}
\label{sec:batch}
The combination of starting population, day of collection, well and cell type, a stratum, gives
cells that are very close in biology, so differences between libraries within a stratum are
mostly technical. We used strata with at least 300 cells from at least two libraries (17 for LSK,
24 for LK). We computed 15 nearest neighbours in our PCA space, which is used throughout the
experiment and is described in Section~\ref{sec:latent}. For each stratum we compared the share of
neighbours from the same library with the share expected under random mixing. A ratio of 1 means
no library structure. The median ratio over strata is 1.03 for LSK (range 0.91--1.09, almost
perfectly mixed) and 1.45 for LK (range 1.07--1.79, a batch effect is present).

On top of these strata we computed an ANOVA of the principal components with library as the
factor (variance explained by library, $\eta^2$, averaged over components with weights equal to
their variance). Library explains 0.2\% of the variance for LSK and 4.6\% for LK.

The LARRY data also let us measure the distance in this space between sister cells of the
same day, well and cell type that were sequenced in different libraries. Their biology should be
as close as possible, so the extra distance compared with sisters from the same library is
attributed to the batch. As a reference, each sister pair was matched with a random pair. We kept
one cell of the pair and replaced its sister with a random cell from another clone in the same
stratum. If the sisters came from one library, the random cell was taken from that library too.
If they came from different libraries, the random cell was taken from a library other than that
of the kept cell. In PCA space, sisters from different libraries are further apart than sisters
from the same library by 0 to 0.5 in median distance, depending on population and day (the
largest difference is for LK at day 4). The median distance between sisters is 5--10, and
between a cell and the random cell 9--14.

We conclude that there is no critical batch effect for the LSK population, and for LK it is
intermediate and we must accept it as a limitation when we will interpret the result.

\section{Design checks}
\label{sec:checks}
All checks in this section use the development clones (1{,}760 clones, 26{,}106 sister pairs in
the two branches). The reported acc is the equal-weight mean over the neutrophil and monocyte
branches; intervals are 95\% bootstrap intervals over clones.

\subsection{Sensitivity to the root}
Five roots were drawn at random from the 100 undifferentiated day-2 cells closest to the centroid
of this population in PCA space. Table~\ref{tab:roots} shows acc for each root on PCA with
$k=15$. The standard deviation over roots is 0.008--0.012, and for Palantir acc ranges from 0.870
to 0.902 depending on the root. This is smaller than the bootstrap interval of a single method
(width about 0.06), but latent spaces are compared on the same clones, so most of the clone
sampling noise cancels in their differences while the effect of the root does not. We therefore
report acc averaged over the five roots for Euclidean distance, DPT and Palantir. Slingshot uses
only the cluster that contains the root, so it is run once.

\begin{table*}[t]
\centering\small
\caption{Accuracy on PCA ($k=15$) for five roots.}
\label{tab:roots}
\setlength{\tabcolsep}{3.5pt}
\begin{tabular}{lcccccccc}
\toprule
Method & 68853 & 16220 & 57472 & 48236 & 60816 & Mean & SD & CI width \\
\midrule
Euclidean distance       & 0.823 & 0.826 & 0.830 & 0.838 & 0.815 & 0.826 & 0.008 & 0.066 \\
DPT ($n_{\mathrm{dc}}=10$) & 0.870 & 0.873 & 0.878 & 0.891 & 0.865 & 0.876 & 0.010 & 0.064 \\
Palantir (20 comp.)      & 0.881 & 0.891 & 0.895 & 0.902 & 0.870 & 0.888 & 0.012 & 0.061 \\
\bottomrule
\end{tabular}
\end{table*}

\subsection{Settings of the trajectory methods}
Settings were compared on PCA with $k=15$ (Table~\ref{tab:settings}; mean over the five roots,
interval averaged over roots). Differences between settings are at most 0.013, well within the
bootstrap interval. DPT uses 10 diffusion components (the Scanpy default). Palantir uses 20
components; this value is only an upper bound for its eigengap-based selection.
\todo{Slingshot}

\begin{table*}[t]
\centering\small
\caption{Accuracy on PCA ($k=15$) for different settings of the trajectory methods.}
\label{tab:settings}
\begin{tabular}{llcc}
\toprule
Method & Setting & acc & 95\% CI \\
\midrule
Euclidean distance & --- & 0.826 & 0.795--0.861 \\
\midrule
DPT & 5 diffusion components  & 0.869 & 0.838--0.903 \\
    & 10 diffusion components & 0.876 & 0.845--0.910 \\
    & 15 diffusion components & 0.875 & 0.844--0.909 \\
    & 20 diffusion components & 0.873 & 0.843--0.908 \\
    & 30 diffusion components & 0.870 & 0.839--0.904 \\
\midrule
Palantir & 5 components  & 0.875 & 0.849--0.901 \\
         & 10 components & 0.882 & 0.853--0.916 \\
         & 20 components & 0.888 & 0.860--0.920 \\
\bottomrule
\end{tabular}
\end{table*}

\subsection{Separate k for frontend and backend graphs}
In the main grid the graph inside a frontend and the graph of the trajectory method use the same
$k$. To check this choice, smoothed PCA ($t=3$) was built with $k_f \in \{10, 15, 30\}$ and DPT
was run with $k_b \in \{10, 15, 30\}$ (root 68853; Table~\ref{tab:kcross}). Equal $k$ is never
worse than unequal $k$ beyond the bootstrap interval. The only large difference is $k_f=30$,
$k_b=10$, where acc drops to 0.701 (0.556 in the neutrophil branch).

\begin{table}[ht]
\centering\small
\caption{Accuracy of DPT on smoothed PCA ($t=3$) for frontend graph size $k_f$ (rows) and
backend graph size $k_b$ (columns). Intervals are about $\pm 0.03$ except for $k_f=30$, $k_b=10$
(0.659--0.756).}
\label{tab:kcross}
\begin{tabular}{lccc}
\toprule
$k_f$ \textbackslash{} $k_b$ & 10 & 15 & 30 \\
\midrule
10 & \textbf{0.862} & 0.864 & 0.863 \\
15 & 0.855 & \textbf{0.866} & 0.865 \\
30 & 0.701 & 0.859 & \textbf{0.866} \\
\bottomrule
\end{tabular}
\end{table}

\subsection{Graph kernel of DPT}
Scanpy weights graph edges with the UMAP kernel by default, while the original DPT uses a Gaussian
kernel with a per-cell width $\sigma_i$ (square root of the median squared distance to the $k$
neighbours). DPT on PCA was rerun with the Gaussian kernel and $\sigma_i$ scaled by 0.5, 1 and 2,
for $k=15$ and $k=30$ (Table~\ref{tab:kernel}; five roots, paired clone bootstrap against the UMAP
kernel with $k=15$). The Gaussian kernel is slightly better (up to $+0.006$), the width matters
little, and all settings stay below Palantir on PCA (0.888). The UMAP kernel is kept in the main
grid.

\begin{table*}[t]
\centering\small
\caption{Accuracy of DPT on PCA for different graph kernels. Difference: paired difference to the
UMAP kernel with $k=15$, with 95\% clone bootstrap interval.}
\label{tab:kernel}
\begin{tabular}{lcccc}
\toprule
Kernel & $k$ & acc & Difference & $p$ \\
\midrule
UMAP & 15 & 0.876 & --- & --- \\
UMAP & 30 & 0.873 & $-0.003$ ($-0.006$, $-0.001$) & 0.005 \\
Gaussian, $\sigma \times 0.5$ & 15 & 0.878 & $+0.002$ ($-0.001$, $+0.005$) & 0.20 \\
Gaussian, $\sigma \times 1$ & 15 & 0.880 & $+0.004$ ($+0.001$, $+0.007$) & 0.008 \\
Gaussian, $\sigma \times 2$ & 15 & 0.880 & $+0.004$ ($+0.001$, $+0.007$) & 0.002 \\
Gaussian, $\sigma \times 1$ & 30 & 0.882 & $+0.006$ ($+0.002$, $+0.009$) & $<0.001$ \\
\bottomrule
\end{tabular}
\end{table*}


\section{Additional results}

\begin{table*}[t]
\centering\scriptsize
\caption{All latent spaces at full depth ($k=15$, five roots, seed 0). PCA: acc; other rows:
paired difference to PCA with 95\% clone bootstrap interval.}
\label{tab:fullall}
\setlength{\tabcolsep}{3pt}
\begin{tabular}{llll}
\toprule
Latent space & Euclidean distance & DPT & Palantir \\
\midrule
PCA-50 & 0.826 & 0.876 & 0.888 \\ \midrule
PCA-5 & $+$0.042 ($+$0.031, $+$0.056) & $-$0.035 ($-$0.053, $-$0.022) & $-$0.016 ($-$0.022, $-$0.011) \\
PCA-10 & $+$0.026 ($+$0.019, $+$0.036) & $-$0.002 ($-$0.006, $+$0.002) & $-$0.019 ($-$0.026, $-$0.014) \\
PCA-20 & $+$0.016 ($+$0.012, $+$0.022) & $-$0.002 ($-$0.008, $+$0.002) & $-$0.016 ($-$0.026, $-$0.009) \\
Smoothed PCA, $t=1$ & $+$0.017 ($+$0.013, $+$0.022) & $-$0.002 ($-$0.007, $+$0.002) & $-$0.002 ($-$0.004, $+$0.001) \\
Smoothed PCA, $t=3$ & $+$0.028 ($+$0.021, $+$0.036) & $-$0.000 ($-$0.004, $+$0.003) & $-$0.006 ($-$0.011, $-$0.002) \\
Smoothed PCA, $t=10$ & $+$0.040 ($+$0.030, $+$0.053) & $-$0.270 ($-$0.350, $-$0.197) & $-$0.007 ($-$0.011, $-$0.004) \\
scVI (default) & $-$0.024 ($-$0.039, $-$0.011) & $-$0.004 ($-$0.011, $+$0.004) & $-$0.021 ($-$0.029, $-$0.013) \\
scVI, 10 dim., KL 0.01 & $+$0.008 ($-$0.006, $+$0.022) & $+$0.003 ($-$0.005, $+$0.010) & $-$0.013 ($-$0.020, $-$0.006) \\
scVI, 10 dim., KL 1 & $-$0.057 ($-$0.073, $-$0.043) & $+$0.008 ($+$0.001, $+$0.016) & $-$0.028 ($-$0.040, $-$0.017) \\
scVI, 30 dim., KL 0.01 & $-$0.023 ($-$0.048, $-$0.001) & $-$0.000 ($-$0.011, $+$0.009) & $-$0.012 ($-$0.023, $-$0.002) \\
scVI, 30 dim., KL 1 & $-$0.047 ($-$0.063, $-$0.034) & $+$0.004 ($-$0.005, $+$0.012) & $-$0.027 ($-$0.041, $-$0.016) \\
scTour & $-$0.012 ($-$0.036, $+$0.011) & $-$0.055 ($-$0.080, $-$0.034) & $-$0.063 ($-$0.085, $-$0.044) \\
GAE-10 & $-$0.217 ($-$0.298, $-$0.146) & $+$0.007 ($-$0.000, $+$0.014) & $-$0.022 ($-$0.038, $-$0.008) \\
GAE-50 & $-$0.212 ($-$0.267, $-$0.164) & $+$0.014 ($+$0.008, $+$0.020) & $-$0.000 ($-$0.003, $+$0.004) \\
GAE-50, untrained & $-$0.050 ($-$0.072, $-$0.032) & $+$0.002 ($-$0.005, $+$0.009) & $-$0.024 ($-$0.033, $-$0.016) \\
VGAE-10 & $-$0.200 ($-$0.276, $-$0.134) & $-$0.006 ($-$0.014, $-$0.001) & $-$0.009 ($-$0.014, $-$0.005) \\
VGAE-50 & $-$0.198 ($-$0.253, $-$0.147) & $+$0.012 ($+$0.007, $+$0.017) & $-$0.000 ($-$0.004, $+$0.003) \\
MAGIC & \todo{} & \todo{} & \todo{} \\
\bottomrule
\end{tabular}
\end{table*}

\begin{table*}[t]
\centering\small
\caption{Accuracy for three training seeds at full depth ($k=15$, five roots).}
\label{tab:seeds}
\begin{tabular}{lcccc}
\toprule
Latent space & Seed & Euclidean distance & DPT & Palantir \\
\midrule
scVI & 0 & 0.802 & 0.872 & 0.867 \\
scVI & 1 & 0.778 & 0.884 & 0.870 \\
scVI & 2 & 0.807 & 0.882 & 0.873 \\
\midrule
GAE-50 & 0 & 0.615 & 0.889 & 0.888 \\
GAE-50 & 1 & 0.602 & 0.889 & 0.888 \\
GAE-50 & 2 & 0.603 & 0.890 & 0.888 \\
\midrule
PCA-50 & --- & 0.826 & 0.876 & 0.888 \\
\bottomrule
\end{tabular}
\end{table*}

\begin{table*}[t]
\centering\scriptsize
\caption{Accuracy after thinning, all depths ($k=15$). PCA: acc; other rows: paired difference to
PCA on the same draw. Mean over draws with minimum and maximum in brackets (three draws for
$p \le 0.5$, one draw otherwise and for GAE-50). At $p=0.1$ the third draw of scVI has 3{,}735
cells with identical latent coordinates and the DPT graph splits into many components.}
\label{tab:thinall}
\setlength{\tabcolsep}{3pt}
\begin{tabular}{lllll}
\toprule
Latent space & $p$ & Euclidean distance & DPT & Palantir \\
\midrule
PCA-50 & 1 & 0.826 & 0.876 & 0.888 \\
 & 0.7 & 0.817 & 0.869 & 0.878 \\
 & 0.5 & 0.810 [0.802, 0.814] & 0.873 [0.869, 0.876] & 0.874 [0.870, 0.878] \\
 & 0.3 & 0.779 [0.775, 0.782] & 0.826 [0.821, 0.833] & 0.845 [0.841, 0.851] \\
 & 0.2 & 0.768 [0.763, 0.774] & 0.763 [0.755, 0.774] & 0.780 [0.766, 0.795] \\
 & 0.1 & 0.772 [0.760, 0.778] & 0.527 [0.509, 0.554] & 0.571 [0.527, 0.607] \\
\midrule
Smoothed PCA, $t=1$ & 1 & $+$0.017 & $-$0.002 & $-$0.002 \\
 & 0.7 & $+$0.022 & $+$0.004 & $-$0.002 \\
 & 0.5 & $+$0.020 [$+$0.019, $+$0.022] & $-$0.002 [$-$0.006, $+$0.001] & $-$0.005 [$-$0.011, $+$0.000] \\
 & 0.3 & $+$0.010 [$+$0.009, $+$0.011] & $-$0.010 [$-$0.014, $-$0.006] & $-$0.018 [$-$0.027, $-$0.010] \\
 & 0.2 & $+$0.004 [$+$0.001, $+$0.006] & $-$0.024 [$-$0.033, $-$0.014] & $-$0.050 [$-$0.069, $-$0.021] \\
 & 0.1 & $-$0.001 [$-$0.003, $+$0.003] & $-$0.042 [$-$0.064, $-$0.008] & $-$0.133 [$-$0.161, $-$0.112] \\
\midrule
Smoothed PCA, $t=3$ & 1 & $+$0.028 & $-$0.000 & $-$0.006 \\
 & 0.7 & $+$0.034 & $-$0.003 & $-$0.004 \\
 & 0.5 & $+$0.036 [$+$0.033, $+$0.038] & $-$0.001 [$-$0.004, $+$0.003] & $-$0.005 [$-$0.010, $-$0.001] \\
 & 0.3 & $+$0.015 [$+$0.012, $+$0.016] & $-$0.003 [$-$0.007, $+$0.004] & $-$0.015 [$-$0.018, $-$0.010] \\
 & 0.2 & $+$0.002 [$-$0.001, $+$0.003] & $-$0.002 [$-$0.018, $+$0.009] & $-$0.050 [$-$0.101, $+$0.004] \\
 & 0.1 & $-$0.006 [$-$0.008, $-$0.002] & $-$0.040 [$-$0.072, $-$0.019] & $-$0.096 [$-$0.123, $-$0.060] \\
\midrule
scVI & 1 & $-$0.024 & $-$0.004 & $-$0.021 \\
 & 0.7 & $-$0.047 & $+$0.021 & $-$0.002 \\
 & 0.5 & $-$0.047 [$-$0.056, $-$0.041] & $+$0.004 [$+$0.002, $+$0.006] & $-$0.017 [$-$0.024, $-$0.014] \\
 & 0.3 & $-$0.021 [$-$0.034, $-$0.010] & $+$0.049 [$+$0.034, $+$0.064] & $+$0.030 [$+$0.018, $+$0.039] \\
 & 0.2 & $-$0.041 [$-$0.051, $-$0.030] & $+$0.094 [$+$0.087, $+$0.098] & $+$0.065 [$+$0.025, $+$0.091] \\
 & 0.1 & $-$0.053 [$-$0.066, $-$0.031] & $+$0.166 [$-$0.007, $+$0.255] & $+$0.131 [$-$0.111, $+$0.296] \\
\midrule
GAE-50 & 1 & $-$0.212 & $+$0.014 & $-$0.000 \\
 & 0.7 & $-$0.257 & $+$0.007 & $-$0.004 \\
 & 0.5 & $-$0.234 & $-$0.001 & $-$0.004 \\
 & 0.3 & $-$0.167 & $-$0.033 & $-$0.066 \\
 & 0.2 & $-$0.177 & $-$0.116 & $-$0.207 \\
 & 0.1 & $-$0.184 & $-$0.097 & $-$0.148 \\
\bottomrule
\end{tabular}
\end{table*}


\begin{figure*}[t]
\centering
\includegraphics[width=\linewidth]{../pics/figS_depth_acc_all.pdf}
\caption{Accuracy by groups of sister pairs for all trajectory methods (columns), three latent
spaces (rows) and all depths (colour). The $x$ axis is the median number of molecules in the
shallower cell of the pairs of a group at the given depth, so the same groups shift to the left as
molecules are removed.}
\label{fig:depthall}
\end{figure*}

\end{document}
